% ID: SEL-D0068-OA-E02
% ESTADO: validada. S068.
\textbf{a)} $f'(x)=3(x-1)(x-3)$. El signo es positivo en $(-\infty,1)$ y $(3,\infty)$, y negativo en $(1,3)$. La función crece donde el signo es positivo y decrece donde es negativo. Hay un máximo local en $x=1$ y un mínimo local en $x=3$.

\textbf{b)} $f''(x)=6(x-2)$. La gráfica es cóncava hacia abajo en $(-\infty,2)$ y convexa, o cóncava hacia arriba, en $(2,\infty)$.

\textbf{c)} $f'(2)=-3$. La recta que pasa por $(2,5)$ con esa pendiente es $y-5=-3(x-2)$, es decir, $\boxed{y=-3x+11}$.
